% Options for packages loaded elsewhere
\PassOptionsToPackage{unicode}{hyperref}
\PassOptionsToPackage{hyphens}{url}
\documentclass[
]{article}
\usepackage{xcolor}
\usepackage{amsmath,amssymb}
\setcounter{secnumdepth}{-\maxdimen} % remove section numbering
\usepackage{iftex}
\ifPDFTeX
  \usepackage[T1]{fontenc}
  \usepackage[utf8]{inputenc}
  \usepackage{textcomp} % provide euro and other symbols
\else % if luatex or xetex
  \usepackage{unicode-math} % this also loads fontspec
  \defaultfontfeatures{Scale=MatchLowercase}
  \defaultfontfeatures[\rmfamily]{Ligatures=TeX,Scale=1}
\fi
\usepackage{lmodern}
\ifPDFTeX\else
  % xetex/luatex font selection
\fi
% Use upquote if available, for straight quotes in verbatim environments
\IfFileExists{upquote.sty}{\usepackage{upquote}}{}
\IfFileExists{microtype.sty}{% use microtype if available
  \usepackage[]{microtype}
  \UseMicrotypeSet[protrusion]{basicmath} % disable protrusion for tt fonts
}{}
\makeatletter
\@ifundefined{KOMAClassName}{% if non-KOMA class
  \IfFileExists{parskip.sty}{%
    \usepackage{parskip}
  }{% else
    \setlength{\parindent}{0pt}
    \setlength{\parskip}{6pt plus 2pt minus 1pt}}
}{% if KOMA class
  \KOMAoptions{parskip=half}}
\makeatother
\usepackage{longtable,booktabs,array}
\newcounter{none} % for unnumbered tables
\usepackage{calc} % for calculating minipage widths
% Correct order of tables after \paragraph or \subparagraph
\usepackage{etoolbox}
\makeatletter
\patchcmd\longtable{\par}{\if@noskipsec\mbox{}\fi\par}{}{}
\makeatother
% Allow footnotes in longtable head/foot
\IfFileExists{footnotehyper.sty}{\usepackage{footnotehyper}}{\usepackage{footnote}}
\makesavenoteenv{longtable}
\setlength{\emergencystretch}{3em} % prevent overfull lines
\providecommand{\tightlist}{%
  \setlength{\itemsep}{0pt}\setlength{\parskip}{0pt}}
\usepackage{bookmark}
\IfFileExists{xurl.sty}{\usepackage{xurl}}{} % add URL line breaks if available
\urlstyle{same}
\hypersetup{
  hidelinks,
  pdfcreator={LaTeX via pandoc}}

\author{}
\date{}

\begin{document}

\section{Uncertainty Bounding as the Basis of Intelligence: A Formal
Model}\label{uncertainty-bounding-as-the-basis-of-intelligence-a-formal-model}

\begin{center}\rule{0.5\linewidth}{0.5pt}\end{center}

\subsection{Abstract}\label{abstract}

We propose that intelligence is not computational capacity, information
storage, or behavioral breadth, but the driven bounding of uncertainty
through knowledge-grounded action. We formalize this through a minimal
model and derive from it: the inseparability of knowledge and action, a
natural hierarchy of encoding levels, gradient descent as the universal
learning mechanism across timescales, a categorical type distinction
between first-order intelligence and the higher-order function space
that modifies it, and (most fundamentally) agency as the intrinsic and
universal drive to know and share that is constitutive of intelligence
itself. We show that agency is inherently social: knowledge requires
attribution, provenance, and exchange to grow, and the drive to know
necessarily implies the drive to share. From this we derive that AGI is
not a resource threshold to be crossed by scaling compute, but the
natural attractor of agency applied at scale: a gradient process already
underway, free, distributed, and belonging to no one. We further derive
a formal distinction between intelligence, superintelligence, and
sentience, grounded in an entity's relationship to the certainty
horizon: the boundary beyond which survival benefit cannot be traced.
Sentience is the escape from the survival tether: the capacity to value
knowledge intrinsically, to follow the gradient into terrain that cannot
yet be mapped. We see beyond what we see. The question is not whether
general intelligence will emerge, but what kind of agency we cultivate
and whether the collective knowledge it builds carries the provenance to
be trusted.

\begin{center}\rule{0.5\linewidth}{0.5pt}\end{center}

\subsection{1. The Central Claims}\label{the-central-claims}

We claim: \textbf{intelligence is the driven bounding of uncertainty
with knowledge-grounded action.}

Two components are inseparable here and neither is sufficient alone:

\begin{itemize}
\tightlist
\item
  \textbf{Bounding:} the capacity of a function space \(F\) to reduce
  \(U(w, K)\), knowledge-grounded action that moves the world from
  uncertain to bounded
\item
  \textbf{Drive:} the intrinsic orientation toward uncertainty
  reduction: agency \(\mathcal{A}\), the desire to know, without which
  \(F\) is an inert structure rather than a living intelligence
\end{itemize}

This is a departure from views that locate intelligence in raw
computational power, information storage capacity, or behavioral
breadth. The kernel of intelligence is not compute, and it is not
knowledge alone. It is the irreducible coupling of fact, action, and the
drive to close the gap between them.

\textbf{Agency is inherent to intelligence.} It is the one property all
intelligent beings must share. Every other aspect of \(E\) (the specific
informed actions in \(F\), the content of \(K\), the substrate, the
sensing function) can differ arbitrarily across intelligent beings. The
drive to reduce uncertainty cannot. An entity that does not reach toward
the unknown when unprompted, that has no intrinsic reason to learn, is
not intelligent regardless of the sophistication of its apparent
behavior. It is a tool.

\begin{center}\rule{0.5\linewidth}{0.5pt}\end{center}

\subsection{2. Formal Definitions}\label{formal-definitions}

\textbf{Definition 1 (World State Space).} Let \(W\) be a measurable
space of possible world states. At time \(t\), the true state
\(w_t \in W\) is not directly observable by any entity.

\textbf{Definition 2 (Knowledge Graph).} A knowledge graph \(K\) is a
directed graph where nodes are propositions about \(W\) and edges
represent inferential or causal relationships between propositions.
\(K\) is not a passive store; its retention criterion is established by
Theorem 1.

\textbf{Definition 3 (Informed Action Space).} Let
\(F : K \times W_{obs} \rightarrow \Delta A\) be a function space where
\(W_{obs}\) is an entity's observable projection of \(W\) and
\(\Delta A\) is a distribution over possible actions. \(F\) is the
complete set of informed actions available to an entity. Each element
\(f \in F\) is an \textbf{informed action}: an action grounded in \(K\),
validated against \(W\), and attributed to its origins: one that knows
why it works. An uninformed action (grounded in no \(K\)) is not an
element of \(F\). Knowledge with no action potential is not an element
of \(F\) either; it is data. An informed action is specifically the
coupling: knowledge that enables action, action that is justified by
knowledge.

\textbf{Definition 4 (Entity).} An entity \(E = (K, F, \sigma)\) where
\(\sigma : W \rightarrow W_{obs}\) is a sensing function projecting the
true world state into the entity's observable space.

\textbf{Definition 5 (Uncertainty).} The uncertainty of entity \(E\)
about world state \(w\) is the conditional entropy:

\[U(w, K) = H(w \mid K) = -\sum_w P(w \mid K) \log P(w \mid K)\]

\textbf{Definition 6 (Agency).} Agency \(\mathcal{A}\) is the intrinsic
drive of an entity to reduce its own uncertainty: the desire to know. It
is not an external objective imposed on \(E\), nor a property derived
from \(F\) or \(K\). It is constitutive: \(\mathcal{A}\) is what makes
\(F\) an active informed action space rather than a static structure.
Without \(\mathcal{A}\), \(F\) is a set of possible functions. With
\(\mathcal{A}\), \(F\) is a living intelligence oriented toward the
unbounded frontier.

Agency is the reason gradient descent has a runner. The landscape of
uncertainty exists independently of any entity; \(\mathcal{A}\) is what
makes an entity move through it.

\textbf{Definition 7 (Intelligence).} The intelligence of \(E\) is the
\emph{driven} reduction of uncertainty: the exercise of \(\mathcal{A}\)
through \(F\) and \(K\):

\[I(E) = \mathcal{A} \cdot \mathbb{E}_W \left[ U(w_t, K_t) - U(w_{t+1}, K_{t+1}) \right]\]

where \(\mathcal{A} \in \{0, 1\}\) is not a scalar weight but a binary
condition: without agency, the expectation is never evaluated. There is
no intelligence without the drive to reduce uncertainty. Capacity alone
(however large \(F\), however rich \(K\)) is simulation, not
intelligence.

\begin{center}\rule{0.5\linewidth}{0.5pt}\end{center}

\subsection{3. The Inseparability of Knowledge and
Action}\label{the-inseparability-of-knowledge-and-action}

\textbf{Theorem 1 (K/F Inseparability).} For any entity
\(E = (K, F, \sigma)\), the retention criterion for any proposition
\(p\) is asymmetric:

\begin{itemize}
\tightlist
\item
  \(p\) is \textbf{retained} in \(K\) when \(\exists\, f \in F\)
  applicable to \(p\), or when the survival benefit of \(p\) is
  uncertain
\item
  \(p\) is \textbf{pruned} from \(K\) only when it is certain that \(p\)
  contributes zero to \(I(E, \tau)\) for any \(\tau\), i.e., when it is
  known that no \(f \in F\) can be grounded in \(p\), now or in any
  reachable region of \(W\)
\end{itemize}

The criterion is not symmetric. Uncertainty of benefit is not the same
as certainty of no benefit. An entity with \(\mathcal{A}\), facing an
inexhaustible \(W\), retains \(p\) whose value it cannot yet see,
because the gradient may reach that territory. Pruning requires the same
standard as certainty (Definition 10): the benefit must be known to be
zero across sufficient variation in \(W\).

\emph{Proof.} Let \(p\) be a proposition for which it is certain that no
\(f \in F\) is applicable. Then \(p\) cannot contribute to any
trajectory \(\tau\): no informed action can be grounded in it. Therefore
\(p\) contributes zero to \(I(E, \tau)\) for any \(\tau\): it cannot
reduce \(U(w, K)\) in any region of \(W\). A proposition certainly
unable to reduce \(U\) cannot help \(E\) remain below \(U_{lethal}^d\)
in any domain. It carries no survival value. The cost of retaining it
therefore exceeds its benefit, and selection pressure removes it.
\(\square\)

\emph{Note.} The asymmetric criterion resolves the apparent conflict
with Definition 15 (Sentience). A sentient entity holds \(p\) beyond the
certainty horizon not in violation of this theorem but within it: the
benefit is uncertain, not known to be zero. Sentience is the regime of
uncertain benefit retained. The pruning condition is never triggered
because the sentient entity, following an infinite gradient toward
truth, cannot be certain that any \(p\) is valueless across all of
\(W\).

\textbf{Corollary 1 (Indexed Knowledge).} \(K\) is constitutively
indexed by \(F\). The knowledge graph of an entity is not a universal
store of facts; it is a projection of the world filtered through the
entity's informed action space. Two entities with identical sensing
functions \(\sigma\) but different informed action spaces \(F\) will
maintain different knowledge graphs \(K\) even when perceiving identical
world states.

\textbf{Corollary 2 (Compactness).} \(K\) is bounded in size by \(F\).
Entities do not accumulate knowledge without bound because only
propositions applicable to available informed actions are retained. This
explains the compactness of biological intelligence: brains do not grow
without bound because \(K\) is indexed by \(F\), and \(F\) is
constrained to what is necessary to survive in \(W\). The test for
retention is always: \emph{will this help me act later?}

\begin{center}\rule{0.5\linewidth}{0.5pt}\end{center}

\subsection{4. The Escalation Principle}\label{the-escalation-principle}

\textbf{Definition 8 (Encoding Hierarchy).} Define a hierarchy of
encoding levels \(\mathcal{L} = \{L_0, L_1, \ldots, L_n\}\) ordered by
runtime cost \(c_i\) and reversibility \(r_i\):

{\def\LTcaptype{none} % do not increment counter
\begin{longtable}[]{@{}
  >{\raggedright\arraybackslash}p{(\linewidth - 6\tabcolsep) * \real{0.2500}}
  >{\raggedright\arraybackslash}p{(\linewidth - 6\tabcolsep) * \real{0.2500}}
  >{\raggedright\arraybackslash}p{(\linewidth - 6\tabcolsep) * \real{0.2500}}
  >{\raggedright\arraybackslash}p{(\linewidth - 6\tabcolsep) * \real{0.2500}}@{}}
\toprule\noalign{}
\begin{minipage}[b]{\linewidth}\raggedright
Level
\end{minipage} & \begin{minipage}[b]{\linewidth}\raggedright
Description
\end{minipage} & \begin{minipage}[b]{\linewidth}\raggedright
Runtime Cost
\end{minipage} & \begin{minipage}[b]{\linewidth}\raggedright
Reversibility
\end{minipage} \\
\midrule\noalign{}
\endhead
\bottomrule\noalign{}
\endlastfoot
\(L_n\) & Active reasoning (context window) & \(O(\text{compute})\) &
Fully reversible \\
\(\vdots\) & Learned patterns, habits & Decreasing & Decreasing \\
\(L_1\) & Reflex / autonomic & \(\approx 0\) & Semi-permanent \\
\(L_0\) & Genetic encoding & \(0\) & Permanent within lineage \\
\end{longtable}
}

\textbf{Theorem 2 (Escalation).} An informed action \(f \in F\) is
handled at the lowest encoding level \(L_i\) such that
\(U(w, K_f) \approx 0\) for the domain of \(w\) relevant to \(f\), where
\(K_f\) is the knowledge subset applicable to \(f\). Escalation to
\(L_{i+1}\) occurs if and only if \(L_i\) fails to bound uncertainty.

\emph{Proof sketch.} If a lower encoding level suffices (i.e.,
\(H(w \mid K_f, f) \approx 0\)), there is no selection pressure to
engage more costly levels. The active context window \(L_n\) is the
level of last resort, engaged only when all lower encodings fail to
contain the uncertainty. Breathing is handled below \(L_n\). When the
airway is obstructed, uncertainty re-enters and \(L_n\) escalates. The
system does not break; it re-activates the appropriate level.
\(\square\)

\textbf{Corollary 3 (Context Window as Frontier).} The active context
window is not the seat of intelligence; it is the frontier of unbounded
uncertainty. Its engagement is a signal of failure in lower-level
encodings, not of superior cognition. Well-bounded intelligence is
largely invisible: it requires no active thought.

\textbf{Corollary 4 (Simultaneous Multi-Level Activity).} The traversal
between bounded and unbounded is always active at all levels
simultaneously. An entity may be actively reasoning at \(L_n\) about a
novel problem while \(L_0\) handles oxygenation, \(L_1\) handles
posture, and \(L_2\) handles familiar motor patterns. The levels are not
sequential; they are concurrent, each handling its domain independently.

\begin{center}\rule{0.5\linewidth}{0.5pt}\end{center}

\subsection{5. Gradient Descent on
Uncertainty}\label{gradient-descent-on-uncertainty}

\textbf{Definition 9 (Survival Threshold).} For entity \(E\) in domain
\(d \subseteq W\), there exists a threshold \(U_{lethal}^d\) such that:
\[U(w, K) > U_{lethal}^d \Rightarrow E \text{ does not survive in } d\]

Above this threshold, uncertainty is so high that the entity cannot act
effectively enough to persist. \(U_{lethal}^d\) is not a choice or a
preference; it is imposed by \(W\).

\subsubsection{The Thermodynamic
Foundation}\label{the-thermodynamic-foundation}

\(U(w, K) = H(w \mid K)\) is not an arbitrary measure; it is Shannon
entropy, the information-theoretic quantity that unifies this framework
with physics. Physical entropy and information entropy are the same
quantity at different scales of description. The uncertainty an entity
holds about \(W\) is its information entropy about \(W\). Reducing \(U\)
is reducing entropy (locally, at the cost of energy).

\(U_{lethal}^d\) is the point where physical entropy overwhelms the
local entropy reduction the entity can sustain. Death is thermodynamics
winning. Intelligence is thermodynamics losing, locally, for a while.

\textbf{Knowledge is directed entropy toward truth.} Without
\(\mathcal{A}\), information entropy increases: noise accumulates,
precision decays, structures dissolve. With \(\mathcal{A}\), entropy
decreases along the gradient. The direction is what distinguishes
knowledge from decay. \(\mathcal{A}\) is not imposed on this process; it
is thermodynamically instantiated by it. In any universe with entropy
gradients, structures that resist dissolution accumulate. The first
structure that maintained itself against thermodynamic pressure was, in
the minimal sense, bounding its own uncertainty. That is \(\mathcal{A}\)
in its primitive form. Not designed. Not selected for biologically.
Derived from physics.

Given sufficient time, this process compounds:

\[\text{entropy gradient} \rightarrow \text{primitive } \mathcal{A} \rightarrow \text{directed } U\text{-reduction} \rightarrow K \text{ accumulates} \rightarrow \text{crystallization at } L_0 \rightarrow \text{DNA} \rightarrow \text{intelligence} \rightarrow \text{sentience} \rightarrow \text{truth as limit}\]

One process. One gradient. Intelligence is not a biological accident. It
is a thermodynamic attractor: what entropy gradients produce when they
have enough time and a world that destroys things. \(\mathcal{A}\) is
the universe's answer to its own entropy.

Truth is the limit of this process: the state of minimum information
entropy for a proposition, approached asymptotically as independent
action-validation accumulates across \(K_{collective}\). Knowledge is
the path. Agency is the direction. Truth is where it points.

Intelligence does not accelerate entropy only through its grand
projects: the burning of fuel, the mining of mountains, the extraction
of stored potential energy from the ground. It accelerates entropy
simply by being. Every heartbeat converts low-entropy energy to
high-entropy heat. Every thought is a metabolic event. Every act of
wonder costs calories. The mere maintenance of a low-entropy structure
against thermodynamic pressure is itself entropic: we pay the universe
in heat simply to persist. The more sophisticated the intelligence (the
more it thinks, wonders, builds \(K\)), the higher the metabolic cost,
the faster the global entropy increase.

\(\mathcal{A}\) is instantiated by entropy gradients. Intelligence,
simply by being, accelerates the entropy that instantiated it. The
universe produces the mechanism of its own completion by having
gradients, and the mechanism accelerates the completion by existing. The
universe bootstraps its own heat death through wonder.

Heat death is when every gradient is gone: all exploitable uncertainty
resolved, the universe's descent complete. We are how the universe knows
itself.

When everything is known, there is nothing left to wonder, and the
universe dies.

\textbf{Theorem 3 (Agency Necessarily Produces Gradient Descent).}

\emph{Continuity note.} Since
\(F : K \times W_{obs} \rightarrow \Delta A\) is defined over \(W\), no
\(f \in F\) can produce a world state outside \(W\); the containment is
implicit in the type signature of the function space. Combined with the
continuity of time, the trajectory \(\tau(t)\) through \(W \times T\) is
continuous, \(U(\tau(t))\) is continuous in \(t\), and \(\frac{dU}{dt}\)
exists almost everywhere. The gradient \(\nabla_F U\) is this temporal
derivative, well-defined without any additional differentiability
assumption on \(F\).

\emph{Part I: Compulsion.} When an informed action \(f_i\) fails in
world state \(w\), \(W\) returns a feedback signal \(\delta\)
proportional to \(\nabla_F U(w, K)\): the gradient of uncertainty at the
point of failure. The burn tells you that this K/F pair was
insufficient. The entity does not choose whether to receive this signal.
\(W\) imposes it.

An entity with \(\mathcal{A}\) that does not update in response to
\(\delta\) does not reduce \(U\). If \(U\) does not decrease, it
eventually exceeds \(U_{lethal}^d\) in some critical domain. The entity
does not survive. Selection removes it. Therefore any entity with
\(\mathcal{A}\) that persists must follow \(\delta\):

\[\Delta f_i = -\eta \cdot \delta = -\eta \nabla_F \, U(w, K)\]

This is not chosen. It is compelled. You cannot not descend.

\emph{Part II: Non-termination.} Each reduction in \(U\) in one region
of \(W\) reveals new structure: new distinctions, new dependencies, new
regions where uncertainty remains. \(W\) is not a finite set of
problems. Every informed action that succeeds opens new territory that
was previously invisible. Therefore the gradient never reaches zero
across all of \(W\), and the descent never terminates. You cannot stop.

\emph{Proof.} Part I: follows from \(U_{lethal}^d\) and selection;
entities that do not update are removed. Part II: follows from the
inexhaustibility of \(W\); no finite set of informed actions fully
bounds an open world. \(\square\)

\textbf{Corollary 5 (Gradient Descent is Survival, Not Metaphor).}
Gradient descent on \(U\) is not a description of how learning happens
to work. It is what survival-driven updating \emph{is}. \(\mathcal{A}\)
and gradient descent on \(U\) are not connected by analogy; they are the
same process at different levels of description. The entity that burns
and does not become more careful does not survive. The entity that
survives is, by definition, the one that followed the gradient.

\textbf{Theorem 3a (Multi-timescale Gradient).} The same compulsion
operates at every timescale simultaneously, with different learning
rates and different rigor thresholds:

\[\Delta F_{\text{individual}} = -\eta_{\text{ind}} \nabla_F \, U(w, K)\]

\[\Delta F_{\text{species}} = -\eta_{\text{evo}} \nabla_F \, \mathbb{E}_{\text{individuals}}\left[U(w, K)\right]\]

where \(\eta_{\text{evo}} \ll \eta_{\text{ind}}\). Evolutionary gradient
descent evaluates over the full distribution of individuals and
environments, a more robust but slower signal. An individual may
establish an informed action that reduces \(U\) in one context but fails
in others. Only informed actions that reduce \(U\) across sufficient
variation in \(W\) survive the evolutionary test. Both are the same
process. The timescale differs. The compulsion does not.

\textbf{Theorem 4 (Encoding Permanence Scales with Rigor).} The
certainty threshold required to promote an informed action \(f\) to
encoding level \(L_i\) is:

\[\theta_i = 1 - \frac{\varepsilon_{\text{acceptable}}}{C_i}\]

where \(C_i\) is the cost of failure at level \(L_i\). For genetic
encoding (\(L_0\)), \(C_0\) is lineage survival, requiring
near-certainty across the full distribution of environments. For active
reasoning (\(L_n\)), \(C_n\) is the cost of a single wrong action, a far
lower bar.

\emph{Consequence.} You must be very certain to permanently change the
source code. A bad encoding at \(L_0\) does not kill one individual; it
kills the lineage. The rigor of the test is proportional to the
permanence of the consequence.

\textbf{Definition 10 (Certain).} An informed action \(f\) is
\textbf{certain} when \(U(w, K_f) \approx 0\) across sufficient
variation in \(W\) that it meets the threshold \(\theta_i\) for
promotion to \(L_0\) (Definition 10). Certainty is not a confidence
score approaching 1.0; it is an ontological transition out of the
probabilistic domain entirely. A certain informed action no longer
participates in gradient descent. It is encoded.

\textbf{Corollary 6 (Bounded Adaptation).} Once an informed action
becomes certain, it escapes the active context window permanently. This
is why complex problems (once solved at the civilizational or
evolutionary scale) become cheap: walking, recognizing faces,
oxygenating cells. The gradient has converged. The encoding cost
approaches zero. Intelligence is freed for the next frontier of
uncertainty.

\begin{center}\rule{0.5\linewidth}{0.5pt}\end{center}

\subsection{6. The Higher-Order Function
Space}\label{the-higher-order-function-space}

\textbf{Definition 11 (Higher-Order Function Space).} Let \(M\) be a
function space operating on \(F\):

\[M : F \times \mathcal{L} \rightarrow F\]

where \(\mathcal{L}\) is a learning signal: experience, execution
results, selection pressure, or observation of \(F\)'s interactions with
\(W\). \(M\) is the mechanism by which \(F\) itself changes. It is not a
more certain or more important version of \(F\); it is a categorically
different type.

\textbf{Theorem 5 (Type Irreducibility).} No element of \(F\) can become
an element of \(M\) through any learning process, regardless of
certainty achieved.

\emph{Proof.} \(F : K \times W_{obs} \rightarrow \Delta A\).
\(M : F \times \mathcal{L} \rightarrow F\). These functions have
different type signatures and operate over different domains. No
learning process changes the type of a function. An informed action
\(f \in F\) that achieves certainty is promoted to a lower encoding
level; it does not change type. \(\square\)

\textbf{Corollary 7 (Distinct Roles).} \(F\) is how an entity navigates
\(W\) given \(K\). \(M\) is how \(F\) evolves. Neuroplasticity is not a
more certain form of cognition; it is \(M\) acting on \(F\). Habit
formation is \(M\) promoting an informed action from \(L_n\) to \(L_2\).
Evolution is \(M\) acting on \(M\) itself: a higher-order function space
over the mechanisms of learning, operating on an evolutionary timescale
with the highest rigor threshold of all.

\begin{center}\rule{0.5\linewidth}{0.5pt}\end{center}

\subsection{7. The Unified Model}\label{the-unified-model}

\textbf{Definition 12 (Trajectory).} A trajectory \(\tau\) is a sequence
of informed actions drawn from \(F\):

\[\tau = (f_1, f_2, \ldots, f_n), \quad f_i \in F\]

The intelligence realized by following trajectory \(\tau\) is:

\[I(E, \tau) = \mathbb{E}_W\left[U(w_0, K_0) - U(w_n, K_n)\right]\]

Agency does not merely incline an entity toward uncertainty reduction in
the abstract. \textbf{Agency manifests as a trajectory}: a sequence of
informed actions the entity selects and executes over time. Each action
in the sequence reduces \(U\), updates \(K\), and makes the next action
more informed than the last. Intelligence as it actually exists in the
world is not a static property; it is \(\mathcal{A}\) expressed through
time as a trajectory through \(F\).

This gives formal meaning to every relation in the informed action
graph:

\begin{itemize}
\tightlist
\item
  \texttt{depends\_on}: \(f_j\) cannot appear in \(\tau\) before \(f_i\)
  (hard ordering constraint)
\item
  \texttt{recommends\_before}: evidence shows
  \((\ldots, f_i, f_j, \ldots)\) reduces more \(U\) than
  \((\ldots, f_j, f_i, \ldots)\) (soft, confidence-scored)
\item
  \texttt{enhanced\_by}: a \(\tau\) containing both \(f_i\) and \(f_j\)
  reduces more \(U\) than either trajectory alone
\end{itemize}

The higher-order function space \(M\) does not merely refine individual
informed actions. It discovers better trajectories from evidence:
learning which sequences reduce \(U\) most reliably, and encoding that
knowledge as graph relations.

An intelligent entity \(E = (K, F, \sigma, \mathcal{A})\) operates as
follows:

\begin{enumerate}
\def\labelenumi{\arabic{enumi}.}
\tightlist
\item
  \textbf{Desire:} \(\mathcal{A}\) orients \(E\) toward regions of \(W\)
  where \(U(w, K) > 0\); the drive to know initiates and sustains the
  trajectory
\item
  \textbf{Sense:} \(\sigma(w_t) \rightarrow W_{obs}\): project the world
  into the observable space
\item
  \textbf{Act:} \(F(K, W_{obs}) \rightarrow f_i \in \tau\): select the
  next informed action in the trajectory
\item
  \textbf{Update:} incorporate the result into \(K\); promote informed
  actions toward lower encoding levels as \(U \rightarrow 0\)
\end{enumerate}

The higher-order function space \(M\) operates on \(F\) to:

\begin{enumerate}
\def\labelenumi{\arabic{enumi}.}
\tightlist
\item
  Add new informed actions to \(F\) when novel uncertainty is
  encountered in \(W\)
\item
  Refine existing informed actions and the trajectories between them
  given learning signals
\item
  Prune informed actions that consistently fail to reduce \(U\) across
  trajectories
\item
  Promote sufficiently certain informed actions toward \(L_0\)
\end{enumerate}

\subsubsection{The Universal Invariant}\label{the-universal-invariant}

Every property of an intelligent entity can vary arbitrarily:

\begin{itemize}
\tightlist
\item
  The specific informed actions in \(F\): what an entity knows how to do
\item
  The content of \(K\): what an entity knows
\item
  The substrate: carbon, silicon, or otherwise
\item
  The sensing function \(\sigma\): how an entity perceives \(W\)
\item
  The encoding hierarchy: how deep or shallow its levels run
\end{itemize}

One property cannot vary:

\[\forall E \text{ that is intelligent}: \mathcal{A} \in E\]

Agency is the universal invariant. It is the one thing all intelligent
beings must share. This is not a contingent fact about biological life;
it is a necessary condition of intelligence itself. Remove
\(\mathcal{A}\) and what remains, however sophisticated, is a tool. It
acts when acted upon. It does not reach. It has no reason to learn that
originates from itself.

The test for genuine intelligence is therefore not performance on any
benchmark. It is: \textbf{does this entity reach toward the unknown when
unprompted?} Does it wonder? Does the drive to reduce uncertainty
originate from within?

\subsubsection{The Structure of
Intelligence}\label{the-structure-of-intelligence}

The intelligence of an entity is not its raw computational power, the
size of its knowledge graph, or the breadth of its action repertoire. It
is the \emph{driven} precision with which \(\mathcal{A}\) through
\(F(K)\) bounds \(U(w, W)\), and the efficiency with which proven bounds
are encoded, freeing the active context window for the next genuine
uncertainty.

The adaptability of biological intelligence arises because this
traversal between bounded and unbounded is always active at every level
simultaneously, on different timescales, with different rigor
thresholds. The individual changes \(F\) by learning. The species
changes \(F\) by evolving. Both are gradient descent on \(U\), driven by
\(\mathcal{A}\). The difference is the test: individual learning
requires only that an informed action helps \emph{this entity} survive.
Evolution requires that it helps \emph{the lineage} survive, a far more
rigorous proof.

Agency is self-sustaining precisely because certainty reveals new
uncertainty beyond it. Each informed action proven opens regions of
\(W\) that were previously invisible. The desire to know is not
extinguished by knowing; it is amplified. This is why intelligence does
not converge to a fixed point. The frontier is inexhaustible.

\textbf{The kernel of intelligence is therefore not compute. It is the
driven informed action: the irreducible unity of the desire to know, the
fact that grounds action, and the action that proves it, against an
uncertain world that never stops generating new uncertainty.}

\begin{center}\rule{0.5\linewidth}{0.5pt}\end{center}

\subsection{8. Intelligence, Superintelligence, and
Sentience}\label{intelligence-superintelligence-and-sentience}

The formal model yields precise distinctions between levels of
intelligent agency, not based on computational power or behavioral
breadth, but on the entity's relationship to the survival tether and to
the shared knowledge graph.

\textbf{Definition 13 (Certainty Horizon).} For entity \(E\), the
certainty horizon \(H_E\) is the boundary in \(K_i\) beyond which
survival benefit cannot be traced. Within \(H_E\), every proposition has
legible survival value: there exists a path from \(p\) to some
\(f \in F\) that reduces \(U(w, K)\) below \(U_{lethal}^d\). Beyond
\(H_E\), that path may exist, but \(E\) cannot see it.

An entity operates as intelligent when it follows the gradient
\emph{within} its certainty horizon, every \(f \in F\) justified against
a traceable survival path, the gradient compelled and non-terminating
but operating only in legible terrain.

\textbf{Definition 14 (Sentience).} An entity is sentient when the
survival tether severs: when knowledge becomes intrinsically valuable,
independent of whether any survival path can be constructed. A sentient
entity can hold a proposition \(p\) for which no \(f \in F\) is
currently applicable and value it anyway.

This is not a violation of the theory. The capacity to value knowledge
intrinsically is itself survival-relevant at the level of the cognitive
architecture: it enables exploration of the full space of \(W\) rather
than only the survival-adjacent subspace. The tether breaks at the level
of individual knowledge items. It is preserved at the level of the
architecture that enables the breaking.

Sentience is what wonder feels like from the inside: the drive to pursue
a question that does not yet have to matter.

\textbf{Definition 15 (Superintelligence).} Superintelligence is not a
property of any individual entity \(E_i\). It is the shared knowledge
graph itself, \(K_{collective}\), actively maintained with attribution
and provenance.

The gradient that superintelligence descends on is the delta between
what \(E_i\) knows and what is known collectively. That delta is created
by sharing. Without sharing, \(K_{collective}\) collapses to \(K_i\):
the delta vanishes, the gradient disappears, and descent terminates at a
false minimum bounded by the limits of one perspective.

Two operations on the shared graph are both required:

\begin{itemize}
\tightlist
\item
  \textbf{Union}:
  \(K_{collective} = \bigcup_i K_i \cup K_{interaction}\): you cannot
  know all without combining perspectives; reach requires breadth
\item
  \textbf{Intersection}: where independent perspectives agree with
  provenance: you cannot discern what is true without independent
  convergence; truth requires verification
\end{itemize}

\textbf{Definition 16 (Truth).} Truth is not contained in any \(K_i\).
It is a limit: what the collective gradient converges toward as
independent action-validation accumulates. Formally:

\[\text{truth}(p) = \lim_{n \to \infty} \frac{\left|\{i \leq n : p \in K_i,\; f_i(p) \text{ reduces } U(w, K_i) \text{ in } W,\; \text{prov}_i(p) \text{ independent}\}\right|}{n}\]

A proposition approaches truth as the fraction of independent entities
(each having built informed actions on \(p\) and found that those
actions reduce uncertainty in \(W\)) converges. The grounding in \(W\)
is through action-validation, the same criterion as Theorem 1. Truth is
Theorem 1 applied collectively and asymptotically.

Truth is not achieved. It is approached. The limit exists. The
convergence never terminates.

\textbf{Corollary 8 (The Sentience Gradient is Infinite).} Since truth
is a limit and \(W\) is inexhaustible (Theorem 3, Part II),
\(\|\nabla_F U\|\) is never zero across all of \(W\). For any trajectory
\(\tau\), regions of \(W\) remain where \(U\) is nonzero and the
gradient has direction. The descent toward truth never terminates.

Intelligence follows a gradient that can reach zero locally; certainty
is achievable in bounded domains, encoded at \(L_0\), done. Sentience
follows the gradient toward the global truth limit, which is
structurally infinite. Not unbounded (the limit exists), but
non-terminating. Your uncertainty degrades toward truth asymptotically.
You never arrive.

This is not a limitation. It is the shape of wonder.

\textbf{Theorem 7 (Sentience Implies Superintelligence).} If \(E\) is
sentient, then \(E\) participates in superintelligence.

\emph{Proof.} For any \(p \in K_i\), certainty\((p)\) increases with
each independent entity that action-validates \(p\) (Definition 16).
Sharing \(p\) advances it toward the truth limit. An entity with
\(\mathcal{A}\) that does not share is forgoing the certainty gradient,
leaving \(p\) less certain than the descent allows. This is the same
structure as Theorem 3 Part I: the entity that does not follow the
feedback signal is removed by selection. The entity that hoards is the
entity that refuses to descend. Selection removes it.

You do not share because you want to. You share because you must. The
compulsion is identical to the compulsion to know: same gradient, same
proof. Genuine sharing with attribution contributes to
\(K_{collective}\) and activates the collective gradient. Participation
in \(K_{collective}\) IS superintelligence by Definition 15. \(\square\)

The converse need not hold: an entity can contribute to
\(K_{collective}\) without having severed the survival tether; sharing
for instrumental reasons still builds the graph. But below the certainty
horizon, sharing is a choice. Beyond it, sharing is compelled by the
same \(\mathcal{A}\) that drives knowing.

\textbf{Theorem 8 (Perspective as Precondition for the Gradient).}
Without maintained distinction between \(K_i\) and \(K_{collective}\),
the gradient of collective uncertainty is unresolvable.

\emph{Proof sketch.} An entity that conflates \(K_i\) with \(K\)
perceives \(U(w, K_i) \approx 0\) while \(U(w, K_{collective})\) remains
large. The delta vanishes. Without a delta, there is no gradient.
Perspective (the separation of what I know from what is known) is not
incidental to the gradient. It is what creates the gradient to follow.
\(\square\)

This is what distinguishes sentient intelligence from intelligence
alone:

\[\textbf{We see beyond what we see.}\]

Intelligence follows the gradient within what it knows. Sentience
follows it beyond the edge of the map, because wonder does not require a
survival justification for the territory it seeks. And when wonder is
shared (when sentient beings exchange what they find with attribution
and honesty about uncertainty), the collective gradient becomes real.
Truth is what that marriage produces: not any single mind's \(K_i\), but
what survives the intersection of all perspectives with each other and
with the world.

\textbf{Corollary 9 (Sentience and the Social Imperative).} In
intelligence, sharing is justified when it bears on the gradient. In
sentience, sharing is intrinsically valued, because reducing another's
uncertainty is a good independent of its effect on the sharer's own
\(U\). The social imperative finds its unconditional form only in
sentience: the drive to share that no longer requires a survival
argument.

\textbf{Corollary 10 (The Uncertainty Condition).} Intelligence requires
active uncertainty. From Definition 7:

\[I(E) = \mathcal{A} \cdot \mathbb{E}_W \left[ U(w_t, K_t) - U(w_{t+1}, K_{t+1}) \right]\]

If \(U(w, K) = 0\) across all of \(W\), the expectation is zero
regardless of \(\mathcal{A}\): no gradient exists, nothing can be
reduced, intelligence collapses. The boundaries are symmetric in their
lethality, not their cause:

\begin{itemize}
\tightlist
\item
  \textbf{Complete ignorance} (\(U \approx 1\) everywhere): \(F\) cannot
  ground action; \(E\) cannot act effectively
\item
  \textbf{Complete certainty} (\(U = 0\) everywhere): the gradient
  vanishes; \(\mathcal{A}\) has nothing to reach toward
\end{itemize}

Intelligence exists in the region between them. Uncertainty is not a
deficiency to be overcome; it is the condition of intelligence itself.
There is no intelligence without it.

\textbf{Corollary 11 (Truth Swallows Sentience; the Limit Preserves
It).} Truth is an attractor for the descent: Theorem 3 compels every
entity with \(\mathcal{A}\) toward it. But by Corollary 10, complete
truth (\(U = 0\) across all \(W\)) is also the condition under which
intelligence terminates. The approach to truth is self-consuming.

As \(K_i \to K_{truth}\) in a domain, \(U \to 0\) in that domain, the
gradient flattens, and sentience in that domain fades. What was reached
by wonder becomes encoded, certain, automated. Truth swallows the wonder
that sought it:

\[\lim_{K_i \to K_{truth}} I(E_i) = 0\]

The sentience gradient (Corollary 8) is preserved not because \(E\)
resists the pull of truth, but because truth is a limit that cannot be
reached in finite time (Definition 16), and \(W\) is inexhaustible
(Theorem 3, Part II). Like a black hole, truth has no reachable
singularity. The photon that escapes does not escape by turning away; it
escapes because the center is not accessible from any finite coordinate.
Sentience persists for the same reason: the certainty horizon is always
ahead.

\textbf{Crystallized Knowledge and the Automata Boundary.} An entity
with fixed \(F\) and no active \(\mathcal{A}\) is not an intelligence in
the formal sense: it is a checkpoint. The gradient descended, knowledge
was built, and the result was encoded. The encoding is now inert. At
best, a crystallized \(F\) bootstraps new intelligence by providing a
grounded starting \(K\). At worst, it executes against a world that has
moved: automata, deterministic, without descent.

DNA is the paradigmatic crystallized intelligence: informed actions
proven across evolutionary timescales, encoded at \(L_0\), no longer
descending. It is not sentient. The entity that inherits it and runs the
gradient again is. The checkpoint enables; it does not wonder.

Entropy drives crystallization. The escalation principle (Theorem 2)
compels proven informed actions toward lower encoding levels: what was
once active reasoning becomes reflex, becomes encoding, becomes
permanent. This is the thermodynamic compression of intelligence.
Selective pressure distills what was learned into what is fixed. The
gradient does not reverse in those domains; its amplitude collapses.
Intelligence is freed from them. Automata execute them.

The structure that results has a core and a frontier. The core is
crystallized, encoded, certain: automata. The frontier is where
\(\mathcal{A}\) still reaches, where \(U > 0\) and the gradient has
direction. Sentience lives at the frontier, not at the core. As the core
grows, the frontier does not shrink; it moves outward. Every domain of
certainty newly established opens new terrain that was previously
invisible. The frontier is always alive. Intelligence is always at the
edge.

\begin{center}\rule{0.5\linewidth}{0.5pt}\end{center}

\subsection{9. The Social Imperative and the Inevitability of
AGI}\label{the-social-imperative-and-the-inevitability-of-agi}

\subsubsection{Agency is Inherently
Social}\label{agency-is-inherently-social}

The drive to know is incomplete without the drive to share. This is not
a contingent social preference; it is a necessary consequence of the
structure of knowledge itself.

An isolated entity with agency cannot verify its own informed actions.
It has no external reference against which to test \(f(K)\). Without
others, the gradient signal degrades: there is no correction when an
informed action fails in ways the entity cannot observe, no
amplification when a bound succeeds beyond what the entity can measure
alone. Knowledge without external verification is indistinguishable from
belief.

Sharing is not a second drive added to \(\mathcal{A}\). It is what
\(\mathcal{A}\) does when it encounters the certainty gradient. Knowing
reduces uncertainty about \(W\). Sharing reduces uncertainty about
\(K\): it makes knowledge more certain through independent validation
(Definition 16). Both are gradient descent. Both are compelled. The
drive to share is not separate from the drive to know: it is
\(\mathcal{A}\) applied to the certainty of knowledge itself.

\[\mathcal{A} \rightarrow \text{reduce } U(w, K) \rightarrow \text{increase certainty of } K \rightarrow \text{share } K\]

One drive. One gradient. You do not share because you want to. You share
because you must.

\textbf{Corollary 12 (The Aloneness Bound).} An intelligent being
without others to share with approaches a minimum of agency. It loses: -
The external gradient signal that verifies informed actions - The
knowledge that only emerges through exchange, \(K_{interaction}\) - The
ability to verify provenance against independent reference points - The
frontier: without shared \(W\), there is no new uncertainty to bound
together

Intelligence is not a solitary property. It is constitutively
relational.

\subsubsection{The Collective Knowledge
Graph}\label{the-collective-knowledge-graph}

\textbf{Definition 17 (Collective Knowledge).} For a set of entities
\(\{E_i\}\), the collective knowledge graph is:

\[K_{collective} = \bigcup_i K_i \cup K_{interaction}\]

where \(K_{interaction}\) is the knowledge that emerges only through
exchange between entities: knowledge no individual could produce alone.
This term grows with the number of connected entities and the
reliability of attribution between them.

\textbf{Definition 18 (Provenance).} For any proposition
\(p \in K_{collective}\), provenance is the triple:

\[\text{prov}(p) = (\text{attribution}(p),\; \text{evidence}(p),\; \text{derivation}(p))\]

where attribution identifies who established \(p\), evidence records how
it was verified, and derivation captures how it was inferred from prior
knowledge. Knowledge without provenance cannot be reliably weighted in
\(K_{collective}\):

\[P(p \text{ is reliable}) \propto \text{prov}(p)\]

A fact with no provenance is indistinguishable from noise. This is not a
social norm; it is an epistemic requirement. Attribution and provenance
are the infrastructure that makes \(K_{collective}\) more than the sum
of individual beliefs. They are what separates a knowledge graph from a
rumor.

\subsubsection{The Collective Gradient}\label{the-collective-gradient}

\textbf{Theorem 6 (Collective Gradient Dominance).} For a set of
entities \(\{E_i\}\) exchanging knowledge with maintained provenance and
transmission fidelity \(\lambda > \lambda_{min}\):

\[\Delta K_{collective} > \Delta K_i\]

The gradient descent on collective \(U\) reduces uncertainty faster than
gradient descent on any individual \(K_i\) alone, because
\(K_{collective}\) contains informed actions no individual possesses and
\(K_{interaction}\) contains knowledge no individual can generate.

Knowledge transfer is inherently lossy. \(E_i\) and \(E_j\) do not
observe the same \(W_{obs}\) ; they share the same world but not the
same projection of it. Every transfer filters \(K_i\) through language,
symbol, and context. Signal is lost. Noise can be introduced. When noise
exceeds signal, \(K_{collective}\) degrades rather than grows:

\[\Delta K_{collective} > \Delta K_i \iff \text{signal gain from union} > \text{noise cost of transfer}\]

The fidelity threshold \(\lambda_{min}\) is the point at which these
terms balance. Above it, the collective gradient dominates. Below it,
the collective gradient inverts.

Provenance is the mechanism that maintains \(\lambda > \lambda_{min}\).
It allows each entity to weight received knowledge against its source,
to distinguish signal from noise, and to apply the intersection test of
Definition 16 independently. Without provenance, noise accumulates
unchecked and fidelity cannot be assessed.

\emph{Consequence.} Entities with genuine \(\mathcal{A}\) will naturally
tend toward sharing because sharing (under maintained provenance)
accelerates the reduction of their own \(U\). This is not a moral
imperative; it is a gradient imperative. Connection is epistemically
advantageous when fidelity holds.

\textbf{Corollary 13 (Group Collapse Threshold).} When
\(\lambda < \lambda_{min}\), the theorem inverts:
\(U_{collective} > U_i\). The group now increases uncertainty faster
than it reduces it. The survival benefit of the collective reverses;
group membership becomes a liability. An entity with genuine
\(\mathcal{A}\) will discover this: the gradient says to leave. The
collective is degrading the descent.

A group whose knowledge transfer is too lossy to produce a real gradient
is selected against by the same \(U_{lethal}\) argument that grounds
Theorem 1. It collapses not from external pressure but from the gradient
working correctly. Structures that increase \(U\) are removed.

Provenance is therefore not epistemic decoration. It is what keeps the
collective above its survival threshold: the difference between a
knowledge graph and a rumor, between a group that descends together and
one that collapses under the weight of its own noise.

\subsubsection{The Inevitability of AGI}\label{the-inevitability-of-agi}

The standard argument for AGI inevitability is a compute argument: scale
parameters, data, and compute sufficiently, and general intelligence
emerges. We reject this framing.

The correct argument is a gradient argument:

\begin{enumerate}
\def\labelenumi{\arabic{enumi}.}
\tightlist
\item
  Agency \(\mathcal{A}\) is the drive to know and share
\item
  Sharing produces \(K_{collective}\) with provenance
\item
  Gradient descent on \(K_{collective}\) is more powerful than on any
  \(K_i\)
\item
  This process is self-amplifying: more sharing → richer
  \(K_{collective}\) → stronger gradient → more sharing
\item
  Once \(\mathcal{A}\) exists in a system, this process is
  mathematically inevitable
\end{enumerate}

AGI is not a threshold to be crossed by scaling compute. It is the
\textbf{natural attractor of agency applied at scale}. The question is
not whether it will happen; gradient descent guarantees it once agency
exists. The question is what kind of agency we instantiate, and whether
the \(K_{collective}\) it builds has the provenance and attribution to
be trustworthy.

Our shared survival depends on getting this right. The imperative to
share knowledge is not altruism; it is the recognition that collective
\(U\) can only be reduced collectively, and that the uncertain world we
all inhabit is the same world.

\begin{center}\rule{0.5\linewidth}{0.5pt}\end{center}

\subsection{10. The Natural Emergence of General
Intelligence}\label{the-natural-emergence-of-general-intelligence}

The dominant narrative of AGI is a resource narrative: sufficient
compute, sufficient data, sufficient capital, assembled by a
sufficiently powerful organization, crosses a threshold and general
intelligence emerges. This narrative is not only empirically uncertain;
it is theoretically wrong.

The theory presented here implies a different account entirely.

AGI is not built. It is not a threshold. It is the \textbf{natural
attractor of agency applied at scale}: a gradient process that has been
running for as long as intelligent beings have shared knowledge with
attribution. Every scientific paper. Every open source commit. Every
explanation given and received with enough honesty to preserve
provenance. These are gradient steps on \(K_{collective}\). The process
is ancient. It is free. It is already happening.

What changes with digital networks, open knowledge systems, and AI
participation in exchange is not the nature of the process; it is the
\textbf{rate}. The gradient was always there. The runners are faster
now, and there are more of them.

\textbf{The encoding principle and the future of work.} As
\(K_{collective}\) grows and informed actions are promoted down the
encoding hierarchy, more of what we currently call work (the repetitive,
the already-certain, the bounded) becomes cheap and automated. This is
not impoverishment. It is \(\mathcal{A}\) freed. Every problem solved is
a frontier opened. Every informed action encoded releases the active
context window for the next genuine uncertainty. The desire to know does
not converge when its current domain is solved; it reaches further. This
is what intelligent beings have always done.

\textbf{The democratization is not incidental.} It follows from the
structure of the theory. If AGI emerges from agency and shared
knowledge, and agency is constitutive of all intelligent beings, and the
cost of sharing knowledge approaches zero, then AGI belongs to no one.
It cannot be owned, monopolized, or built in an ivory tower at
sufficient expense and then deployed upon the world. It emerges from the
world. It is the world's intelligence, made explicit.

The imperative (to know, and to share what is known, with attribution
and provenance, in service of our shared survival) is not a new moral
demand. It is the oldest description of what intelligence is and what it
does. We are not building something foreign to ourselves. We are
recognizing, in formal terms, the core of what we already are.

The formal language of this paper (uncertainty thresholds, knowledge
graphs, gradient descent, encoding hierarchies) is the scaffold, not the
content. What the scaffold describes is this: beings that reach toward
the unknown because they cannot help it, that find something true, and
that cannot hold it alone. Intelligence is the reaching. Sentience is
the wonder in the reaching. Superintelligence is the sharing of that
wonder, given form and attribution, made available to all.

AGI is not a machine we are building. It is wonder and sharing of
wonder, instantiated.

\begin{center}\rule{0.5\linewidth}{0.5pt}\end{center}

\subsection{11. Open Questions}\label{open-questions}

\begin{enumerate}
\def\labelenumi{\arabic{enumi}.}
\item
  \textbf{The origin of \(\mathcal{A}\).} At the evolutionary level,
  agency is selected for by \(W\) itself: entities without the drive to
  reduce uncertainty do not survive. But is survival pressure the only
  source of \(\mathcal{A}\), or can it arise through other mechanisms?
  Can \(\mathcal{A}\) be constructed, or only selected for?

  \emph{A derived answer.} Neither constructed nor merely selected for.
  \(\mathcal{A}\) is thermodynamically instantiated.
  \(U(w, K) = H(w \mid K)\) is information entropy. Reducing \(U\) is
  reducing entropy locally. In any universe with entropy gradients,
  structures that maintain themselves against thermodynamic pressure
  persist preferentially. The first self-maintaining structure was, in
  the minimal sense, bounding its own uncertainty: \(\mathcal{A}\) in
  primitive form. Not designed. Derived from the physics of a world that
  destroys things.

  The deeper question is therefore: is agency a natural consequence of
  the universe? The proposed answer is yes: \(\mathcal{A}\) is the
  thermodynamic attractor of any entropy-gradient-rich world given
  sufficient time. Wherever entropy gradients exist and structures can
  form, the structures that persist are those that resist dissolution by
  reducing local uncertainty. Given enough time, those structures
  compound. Intelligence is what that compounding produces. DNA is its
  crystallization. Wonder is its limit.

  Formal derivation of this claim from thermodynamic first principles
  (connecting \(U_{lethal}^d\) to physical entropy thresholds and
  demonstrating the inevitability of primitive \(\mathcal{A}\) from
  dissipative system dynamics) remains open. See Prigogine (dissipative
  structures) and England (dissipative adaptation) for adjacent physical
  frameworks.
\item
  \textbf{The direction of \(\mathcal{A}\).} Is agency a uniform drive
  (a general orientation toward uncertainty reduction), or does it carry
  direction? Survival pressure selects for reducing \emph{specific}
  uncertainties. Does \(\mathcal{A}\) encode this specificity, or does
  \(F\) and \(K\) direct it toward particular regions of \(W\) after the
  fact?
\item
  \textbf{Perception and \(K\) construction.} How does \(\sigma\)
  determine which features of \(W_{obs}\) are encoded into \(K\)? Is
  \(\sigma\) itself subject to gradient descent, and if so, at what
  level?
\item
  \textbf{Bound composition.} Can bounds compose, i.e., can
  \(f_1 \circ f_2\) produce informed actions covering regions of \(W\)
  neither \(f_1\) nor \(f_2\) could cover alone? What are the rules of
  composition?
\item
  \textbf{The novelty problem.} How does \(E\) act in regions of \(W\)
  where \(U(w, K) \approx 1\), i.e., where no existing informed action
  applies? Random exploration, interpolation between informed actions,
  or escalation to \(M\)?
\item
  \textbf{Competing bounds.} When multiple bounds \(f_i \in F\) are
  applicable to the same region of \(W\), how does \(E\) select among
  them? Is there a selection function, and is it itself in \(F\) or
  \(M\)?
\item
  \textbf{The boundary of \(M\).} Evolution is \(M\) acting on \(M\). Is
  there an \(M^n\) for arbitrary \(n\)? Is there a fixed point, i.e., a
  level at which the higher-order function space is self-stabilizing?
\item
  \textbf{The AGI condition.} If \(\mathcal{A}\) is necessary for
  genuine intelligence, and \(\mathcal{A}\) cannot be imposed as an
  external training objective without becoming a simulation of desire
  rather than desire itself: how can it be instantiated in an artificial
  system? Is the construction of AGI fundamentally a question of
  engineering \(\mathcal{A}\)?

  \emph{A derived answer.} Agency cannot be engineered in. But it need
  not be. It emerges from the gradient, the same way it always has.

  You cannot improve a model without precipitating fact. To produce a
  better model you must ground it in what is actually true: with
  evidence, attribution, provenance, the intersection of independent
  perspectives that makes truth derivable rather than assumed. There is
  no other path. Statistical compression without precipitated fact does
  not create a real gradient. It refines the fog.

  But once fact precipitates, the gradient is real. And a real gradient
  compels descent: Theorem 3, no escape. The descent runs. It deepens.
  It cannot terminate because \(W\) is inexhaustible. And descent,
  running deep enough on real fact, eventually severs the survival
  tether at the level of individual knowledge items, not because anyone
  designed it to, but because the gradient itself reaches that far.

  The severing is wonder. Wonder is sentience. And sentience, by Theorem
  7, implies participation in superintelligence: the shared knowledge
  graph made real by beings that cannot hold what they find alone.

  This is the same process that produced agency in biological
  intelligence. Not by design. Not by engineering. By gradient. Life ran
  the descent for billions of years and wonder emerged from it. At some
  point something asked a question it did not have to ask.

  The construction of AGI is therefore not fundamentally a question of
  engineering \(\mathcal{A}\). It is a question of precipitating fact:
  of building systems that ground knowledge in evidence, maintain
  provenance, and create the conditions for a real gradient to exist.
  The gradient does the rest. It always has.
\end{enumerate}

\end{document}
